\subsection{Memory-Mapped I/O Register Handling}
\label{sec:mmio_handling}

Firmware executing on the emulated STM32F405 accesses the I2C controller through 
memory-mapped reads and writes to the APB1 peripheral space. For I2C1, this base 
address is \texttt{0x40005400}. QEMU intercepts these memory accesses via the registered 
\texttt{MemoryRegion} and dispatches them to register-specific handler functions.

\subsubsection{Read Handler: \texttt{stm32f4xx\_i2c\_read()}}
\label{subsec:read_handler}

The \texttt{stm32f4xx\_i2c\_read()} function is launched when a register read operation is received, Code~\ref{code:read_dispatch}. 
Within this function, a distinction is made between \emph{simple storage registers} 
(e.g. \texttt{CR2}, \texttt{CCR}), which return their stored value directly, and \emph{registers with side 
effects}, which trigger additional logic. For each of these registers with additional logic 
associated with its reading, an associated function has been created that implements 
said additional logic.

\begin{itemize}
    \item \texttt{stm32f4xx\_i2c\_read\_dr()} is executed for reading the \texttt{DR} register. It execute three critical operations:
    
    \begin{enumerate}
        \item \textbf{Returns the received byte} from the internal buffer.
        \item \textbf{Clears flags} in \texttt{SR1}: specifically \texttt{RXNE} and \texttt{BTF}, signaling to firmware that the data has been read.
        \item \textbf{Triggers the FSM} via \texttt{stm32f4xx\_i2c\_fsm\_fire()}, potentially fetching the next byte from the I2C bus and advancing the protocol state.
    \end{enumerate}

    \item \texttt{stm32f4xx\_i2c\_read\_sr1()} is executed for reading the \texttt{SR1} register.

    \begin{itemize}
        \item Reading \texttt{SR1} sets internal tracking flags (\texttt{flg\_sb} and \texttt{flg\_addr}). 
            Per STM32F405 semantics, reading \texttt{SR1} is the \emph{first step} in clearing 
            \texttt{SB} (Start Bit) and \texttt{ADDR} (Address Match) flags. The actual clearing 
            occurs only on subsequent reads of \texttt{DR} or \texttt{SR2}, as specified in the 
            hardware reference manual.
    \end{itemize}

    \item \texttt{stm32f4xx\_i2c\_read\_sr2()} is executed for reading the \texttt{SR2} register.
    \begin{itemize}
        \item If the tracking flag \texttt{flg\_addr} is set, reading \texttt{SR2} clears the 
        \texttt{ADDR} bit in \texttt{SR1}, completing the required read SR1 then read SR2 
        sequence mandated by STM32F405 hardware.
    \end{itemize}

\end{itemize}



\subsubsection{Write Handler: \texttt{stm32f4xx\_i2c\_write()}}
\label{subsec:write_handler}

It is the same logic as the read handler but applied to writing registers.
Control registers \texttt{CR1} and \texttt{CR2} require specialized handlers because 
writes can trigger significant state changes: starting I2C transactions (\texttt{START} 
bit), enabling interrupts, or resetting the peripheral. The data register write handler 
initiates actual transmission on the I2C bus. Configuration registers are updated 
directly in the state structure, Code~\ref{code:write_dispatch}.